%! AP Statistics — Unit 1, Lesson 1.8 — Slides (Beamer)
\documentclass[aspectratio=169, 11pt]{beamer}
\usepackage{apstats-beamer}

%=============================================================================
\begin{document}

% ─────────────────────────────────────────────────────────────────────────────
% SLIDE 1 — LESSON TITLE
% ─────────────────────────────────────────────────────────────────────────────
{
\setbeamercolor{background canvas}{bg=navy}
\begin{frame}[plain]
  \begin{tikzpicture}[remember picture, overlay]
    \fill[navylight] (current page.north west)
      rectangle ([yshift=-1.35cm]current page.north east);
  \end{tikzpicture}
  \vspace{2.8cm}
  \hspace{0.55cm}%
  \begin{minipage}{0.9\textwidth}
    {\color{goldacc}\bfseries\small LESSON 1.8}\\[0.5cm]
    {\color{white}\bfseries\LARGE Graphical Representations of\\[0.25cm]
      Summary Statistics}\\[1.4cm]
    {\color{skymid}\small AP Statistics~$\cdot$~Unit 1: Exploring One-Variable Data}
  \end{minipage}
\end{frame}
}

% ─────────────────────────────────────────────────────────────────────────────
% SLIDE 2 — WARM-UP
% ─────────────────────────────────────────────────────────────────────────────
{
\setbeamercolor{background canvas}{bg=sky}
\begin{frame}[t, plain]
  \navyheader{Warm-Up}
  \vspace{1em}%
  \hspace{0.55cm}%
  \begin{minipage}{0.9\textwidth}
    \small
    \textit{Use the boxplot on your warm-up sheet showing homework hours per week
    for 20 students.}

    \vspace{0.5em}
    \begin{enumerate}
      \setlength{\itemsep}{0.6em}\setlength{\topsep}{0em}
      \item What are the minimum and maximum values? What is the range?
      \item Identify $Q_1$, $M$, and $Q_3$. Calculate the IQR.
      \item Would a student reporting 14 hours be flagged as an outlier?
            Show the upper fence.
      \item Is the distribution symmetric, skewed left, or skewed right?
            How can you tell from the boxplot?
    \end{enumerate}
    \vspace{0.4cm}
    {\color{navy}\itshape\small
      Today we build these from scratch --- and learn what to do when an
      outlier appears.}
  \end{minipage}
\end{frame}
}

% ─────────────────────────────────────────────────────────────────────────────
% SLIDE 3 — PRIMARY OBJECTIVE
% ─────────────────────────────────────────────────────────────────────────────
{
\setbeamercolor{background canvas}{bg=sky}
\begin{frame}[t, plain]
  \begin{tikzpicture}[remember picture, overlay]
    \fill[navy] (current page.north west)
      rectangle ([yshift=-0.25cm]current page.north east);
  \end{tikzpicture}
  \vspace{0.45cm}\par
  \hspace{0.55cm}%
  \begin{minipage}{0.9\textwidth}

    \sectionlabel{Primary Objective}

    \normalsize
    Students will \textbf{read and interpret} boxplots and modified boxplots,
    apply the $1.5 \times \text{IQR}$ outlier rule, and \textbf{describe
    distributions} using SOCS.

    \vspace{0.5cm}
    \textbf{Learning Targets:}
    \begin{enumerate}\small
      \setlength{\itemsep}{0.4em}
      \item Read a boxplot: identify the five-number summary, IQR, and range.
      \item Determine outliers using the fence rule; interpret modified boxplots.
      \item Describe a distribution from a boxplot using SOCS.
      \item Explain what a boxplot can and cannot show versus a histogram.
    \end{enumerate}
  \end{minipage}
\end{frame}
}

% ─────────────────────────────────────────────────────────────────────────────
% SLIDE 4 — HOOK: FIVE NUMBERS → ONE PICTURE
% ─────────────────────────────────────────────────────────────────────────────
{
\setbeamercolor{background canvas}{bg=goldbg}
\begin{frame}[t, plain]
  \navyheader{Hook: Five Numbers $\to$ One Picture}
  \vspace{0.8em}%
  \hspace{0.55cm}%
  \begin{minipage}{0.9\textwidth}
    \small
    A restaurant recorded delivery times (min) for 15 orders.\\
    Five-number summary: $\min=12$,\ $Q_1=20$,\ $M=27$,\ $Q_3=35$,\ $\max=48$

    \vspace{0.5em}
    \begin{center}
    \begin{tikzpicture}[scale=1.1]
      \draw[thick] (0.6,0) -- (6.6,0);
      \foreach \v/\lbl in {0.60/10, 1.20/15, 1.80/20, 2.40/25, 3.00/30,
                            3.60/35, 4.20/40, 4.80/45, 5.40/50} {
        \draw (\v,-0.12) -- (\v,0.12);
        \node[below, font=\small] at (\v,-0.12) {\lbl};
      }
      \node[below, font=\small\itshape] at (3.0,-0.56) {Delivery time (minutes)};
      \draw[thick, navy] (0.84, 0.38) -- (1.80, 0.38);
      \draw[thick, navy] (1.80, 0.10) rectangle (3.60, 0.66);
      \draw[thick, navy] (2.64, 0.10) -- (2.64, 0.66);
      \draw[thick, navy] (3.60, 0.38) -- (5.16, 0.38);
      \draw[thick, navy] (0.84, 0.22) -- (0.84, 0.54);
      \draw[thick, navy] (5.16, 0.22) -- (5.16, 0.54);
    \end{tikzpicture}
    \end{center}

    \vspace{0.3em}
    \begin{itemize}
      \setlength{\itemsep}{0.4em}
      \item What does the \textbf{box} represent? \quad What does each \textbf{whisker} represent?
      \item Each of the four sections contains $\approx$ \textbf{25\%} of the data.
      \item What can you say about the \textbf{shape} from this boxplot?
    \end{itemize}
  \end{minipage}
\end{frame}
}

% ─────────────────────────────────────────────────────────────────────────────
% SLIDE 5 — ANATOMY OF A BOXPLOT
% ─────────────────────────────────────────────────────────────────────────────
{
\setbeamercolor{background canvas}{bg=white}
\begin{frame}[t, plain]
  \begin{tikzpicture}[remember picture, overlay]
    \fill[navy] (current page.north west)
      rectangle ([yshift=-0.25cm]current page.north east);
  \end{tikzpicture}
  \vspace{0.45cm}\par
  \hspace{0.55cm}%
  \begin{minipage}{0.9\textwidth}
    \sectionlabel{Anatomy of a Boxplot}

    \small
    \begin{columns}[T]
    \begin{column}{0.45\textwidth}
      \textbf{Five-number summary mapped to boxplot:}
      \begin{itemize}\setlength{\itemsep}{0.3em}
        \item $\min$ \textrightarrow{} left whisker end
        \item $Q_1$ \textrightarrow{} left edge of box
        \item $M$ \textrightarrow{} line inside box
        \item $Q_3$ \textrightarrow{} right edge of box
        \item $\max$ \textrightarrow{} right whisker end
      \end{itemize}

      \vspace{0.5em}
      \textbf{Each section $\approx$ 25\% of data}\\
      Box spans middle \textbf{50\%}\\
      IQR = $Q_3 - Q_1$
    \end{column}
    \begin{column}{0.50\textwidth}
      \textbf{Skewness clues:}
      \begin{itemize}\setlength{\itemsep}{0.3em}
        \item Right whisker longer \textrightarrow{} skewed right
        \item Left whisker longer \textrightarrow{} skewed left
        \item Median off-center in box \textrightarrow{} skewed toward that tail
        \item Equal whiskers, centered median \textrightarrow{} symmetric
      \end{itemize}
    \end{column}
    \end{columns}
  \end{minipage}
\end{frame}
}

% ─────────────────────────────────────────────────────────────────────────────
% SLIDE 6 — MODIFIED BOXPLOT
% ─────────────────────────────────────────────────────────────────────────────
{
\setbeamercolor{background canvas}{bg=white}
\begin{frame}[t, plain]
  \begin{tikzpicture}[remember picture, overlay]
    \fill[navy] (current page.north west)
      rectangle ([yshift=-0.25cm]current page.north east);
  \end{tikzpicture}
  \vspace{0.45cm}\par
  \hspace{0.55cm}%
  \begin{minipage}{0.9\textwidth}
    \sectionlabel{Modified Boxplot}

    \small
    \begin{columns}[T]
    \begin{column}{0.45\textwidth}
      \textbf{Standard vs.\ Modified:}
      \begin{itemize}\setlength{\itemsep}{0.4em}
        \item \textit{Standard}: whisker goes to $\min$ and $\max$.
        \item \textit{Modified}: whisker ends at last \textbf{non-outlier};
              outliers plotted as dots ($\bullet$).
      \end{itemize}

      \vspace{0.6em}
      \textbf{Fence rule reminder:}
      \begin{itemize}\setlength{\itemsep}{0.3em}
        \item Lower fence $= Q_1 - 1.5\,\text{IQR}$
        \item Upper fence $= Q_3 + 1.5\,\text{IQR}$
        \item Value \textbf{beyond} a fence $\Rightarrow$ outlier
      \end{itemize}

      \vspace{0.4em}
      {\color{redacc}\textbf{Common error:} Ending the whisker at the
      fence value instead of the last non-outlier data point.}
    \end{column}
    \begin{column}{0.50\textwidth}
      \textbf{Example} (delivery times, new max = 62 min):\\
      $Q_1=20$, $Q_3=35$, IQR $=15$\\
      Upper fence $= 35+22.5 = 57.5$\\
      $62 > 57.5$ \textrightarrow{} outlier; last non-outlier $= 48$

      \vspace{0.5em}
      \begin{tikzpicture}[scale=0.88]
        \draw[thick] (0.15,0) -- (6.20,0);
        \foreach \v/\lbl in {0.20/10, 1.16/20, 2.12/30, 3.08/40, 4.04/50, 5.00/60} {
          \draw (\v,-0.10) -- (\v,0.10);
          \node[below, font=\scriptsize] at (\v,-0.10) {\lbl};
        }
        \node[below, font=\scriptsize\itshape] at (3.0,-0.44) {Delivery time (min)};
        \draw[thick, navy] (0.392, 0.32) -- (1.16, 0.32);
        \draw[thick, navy] (1.16, 0.08) rectangle (2.76, 0.56);
        \draw[thick, navy] (1.832, 0.08) -- (1.832, 0.56);
        \draw[thick, navy] (2.76, 0.32) -- (3.848, 0.32);
        \draw[thick, navy] (0.392, 0.18) -- (0.392, 0.46);
        \draw[thick, navy] (3.848, 0.18) -- (3.848, 0.46);
        \filldraw[navy] (5.552, 0.32) circle (0.09);
        \node[above, font=\scriptsize, navy] at (5.552, 0.42) {62};
        \node[above, font=\scriptsize, navy] at (3.848, 0.56) {48};
      \end{tikzpicture}
    \end{column}
    \end{columns}
  \end{minipage}
\end{frame}
}

% ─────────────────────────────────────────────────────────────────────────────
% SLIDE 7 — BOXPLOT vs HISTOGRAM
% ─────────────────────────────────────────────────────────────────────────────
{
\setbeamercolor{background canvas}{bg=sky}
\begin{frame}[t, plain]
  \navyheader{Boxplot vs.\ Histogram}
  \vspace{0.8em}%
  \hspace{0.55cm}%
  \begin{minipage}{0.9\textwidth}
    \small
    \begin{columns}[T]
    \begin{column}{0.45\textwidth}
      \textbf{\textcolor{navy}{Boxplot shows clearly:}}
      \begin{itemize}\setlength{\itemsep}{0.4em}
        \item Center (median)
        \item Spread (IQR, range)
        \item Presence of outliers
        \item Approximate skewness
        \item Easy comparison of groups\\
              (side-by-side boxplots)
      \end{itemize}
    \end{column}
    \begin{column}{0.50\textwidth}
      \textbf{\textcolor{navy}{Boxplot hides:}}
      \begin{itemize}\setlength{\itemsep}{0.4em}
        \item Exact shape (unimodal? bimodal?)
        \item Gaps and clusters
        \item Exact frequencies / counts
        \item Whether data are continuous or discrete
      \end{itemize}

      \vspace{0.4em}
      {\color{navy}\textit{Rule of thumb:}
      Use a boxplot for \textbf{comparing groups}.
      Use a histogram for \textbf{examining shape in detail}.}
    \end{column}
    \end{columns}
  \end{minipage}
\end{frame}
}

% ─────────────────────────────────────────────────────────────────────────────
% SLIDE 8 — GUIDED NOTES CHECK
% ─────────────────────────────────────────────────────────────────────────────
{
\setbeamercolor{background canvas}{bg=goldbg}
\begin{frame}[t, plain]
  \navyheader{Quick Check --- Guided Notes}
  \vspace{0.8em}%
  \hspace{0.55cm}%
  \begin{minipage}{0.9\textwidth}
    \small
    A modified boxplot: $Q_1 = 40$, $M = 52$, $Q_3 = 60$.\\
    Left whisker to 30. Right whisker to 78. Outlier dot at 95.

    \vspace{0.6em}
    Answer each:
    \begin{enumerate}\setlength{\itemsep}{0.5em}
      \item What is the IQR?
      \item What is the upper fence? \quad Is 95 above the fence?
      \item What is the range of this dataset?
      \item What does the right whisker endpoint (78) represent?
    \end{enumerate}

    \vspace{0.5em}
    {\color{navy}\small\textit{Key answers: IQR $= 20$;
    fence $= 90$; $95 > 90$, yes; range $= 95 - 30 = 65$;
    78 is the last non-outlier.}}
  \end{minipage}
\end{frame}
}

% ─────────────────────────────────────────────────────────────────────────────
% SLIDE 9 — GROUP ACTIVITY INTRO
% ─────────────────────────────────────────────────────────────────────────────
{
\setbeamercolor{background canvas}{bg=white}
\begin{frame}[t, plain]
  \begin{tikzpicture}[remember picture, overlay]
    \fill[navy] (current page.north west)
      rectangle ([yshift=-0.25cm]current page.north east);
  \end{tikzpicture}
  \vspace{0.45cm}\par
  \hspace{0.55cm}%
  \begin{minipage}{0.9\textwidth}
    \sectionlabel[greenacc]{Group Activity}

    \small
    \textbf{Pulse rates (bpm) for 17 student athletes (sorted):}\\[0.3em]
    52, 56, 58, 60, 62, 63, 64, 66, \textbf{68}, 70, 72, 73, 75, 78, 80, 84, 110

    \vspace{0.6em}
    \textbf{Part 1} (8 min): Compute the five-number summary.\\
    Find $n$, median position, $Q_1$, $Q_3$, IQR, lower and upper fences.\\
    Is 110 an outlier? Where does the right whisker end?

    \vspace{0.4em}
    \textbf{Part 2} (7 min): Read the modified boxplot on your sheet.\\
    Match your calculations to the display. Answer interpretation questions.

    \vspace{0.4em}
    \textbf{Part 3} (5 min): Write a complete SOCS description.\\
    Compare: what does the boxplot show that a histogram would not, and vice versa?
  \end{minipage}
\end{frame}
}

% ─────────────────────────────────────────────────────────────────────────────
% SLIDE 10 — EXIT TICKET
% ─────────────────────────────────────────────────────────────────────────────
{
\setbeamercolor{background canvas}{bg=sky}
\begin{frame}[t, plain]
  \navyheader{Exit Ticket}
  \vspace{0.8em}%
  \hspace{0.55cm}%
  \begin{minipage}{0.9\textwidth}
    \small
    \textit{Use the modified boxplot on your exit ticket --- delivery times for
    20 courier orders.}

    \vspace{0.5em}
    \begin{enumerate}\setlength{\itemsep}{0.5em}
      \item Read the five-number summary and IQR from the boxplot.
      \item Use the $1.5\times\text{IQR}$ fence rule to verify the outlier at
            60 minutes. Show work.
      \item What percent of orders were delivered in 18 to 34 minutes? What
            interval is that?
      \item Why does the right whisker end at 55, not 60?
    \end{enumerate}

    \vspace{0.4em}
    \textbf{AP-Style:} Describe the distribution of delivery times in context.
    Address shape, outlier(s), center, and spread.
  \end{minipage}
\end{frame}
}

% ─────────────────────────────────────────────────────────────────────────────
% SLIDE 11 — LESSON SUMMARY
% ─────────────────────────────────────────────────────────────────────────────
{
\setbeamercolor{background canvas}{bg=navy}
\begin{frame}[plain]
  \begin{tikzpicture}[remember picture, overlay]
    \fill[navylight] (current page.north west)
      rectangle ([yshift=-1.35cm]current page.north east);
  \end{tikzpicture}
  \vspace{1.5cm}
  \hspace{0.55cm}%
  \begin{minipage}{0.9\textwidth}
    {\color{goldacc}\bfseries\large Lesson 1.8 Summary}\\[0.8cm]
    \small
    \begin{itemize}\setlength{\itemsep}{0.5em}
      \item {\color{white} A \textbf{boxplot} places the five-number summary on a number
            line --- each of four sections holds $\approx$25\% of the data.}
      \item {\color{white} A \textbf{modified boxplot} ends whiskers at the last
            non-outlier; outliers appear as separate dots ($\bullet$).}
      \item {\color{white} The whisker ends at a \textit{data point}, not the fence.}
      \item {\color{white} Boxplots reveal center, spread, and outliers; they hide
            exact shape, gaps, and counts.}
      \item {\color{skymid} \textbf{Next:} Lesson 1.9 --- Comparing Distributions with
            side-by-side boxplots.}
    \end{itemize}
  \end{minipage}
\end{frame}
}

\end{document}
